%% --------------------------------------------------------
\chapter{Перспективи та можливі напрями подальших досліджень}
\label{sec:future_research}
%% --------------------------------------------------------

Теоретичні засади та чисельні методи, викладені у цьому посібнику, закладають основу для вирішення низки актуальних проблем сучасної космології. Головна перевага моделей із локальним зламом первинного спектра потужності (broken power spectrum) полягає в можливості розв'язати протиріччя між ранніми та пізніми спостережними даними без руйнування базових успіхів моделі $\Lambda$CDM на великих масштабах.

Нижче сформульовано ключові етапи та відкриті задачі, які становлять безпосередній інтерес для подальшої наукової розробки.

%% --------------------------------------------------------
\section{Чисельне моделювання та байєсівський аналіз}
%% --------------------------------------------------------

Найближчим технічним завданням у межах розглянутого формалізму є реалізація повної чисельної конвеєризації (pipeline) для статистичної перевірки розглянутих теоретичних моделей:
\begin{enumerate}
  \item \textbf{Модифікація лінійних софтверних комплексів:} Інтеграція нестандартного первинного спектра $P_{\mathcal{R}}(k)$ із параметризованим зламом\\ $(k_\mathrm{break}, \Delta n)$ у загальноприйняті космологічні коди лінійної теорії збурень, такі як \texttt{CAMB} або \texttt{CLASS}.
  \item \textbf{Байєсівський пошук параметрів (MCMC):} Запуск багатопараметричного аналізу за методом Монте-Карло за марковськими ланцюгами (MCMC) за допомогою фреймворків \texttt{Cobaya} або \texttt{CosmoMC}.
  \item \textbf{Узгодження даних Planck та JWST:} Головною метою аналізу є пошук спільної області параметрів у просторі $\{k_\mathrm{break}, \Delta n, \Omega_b h^2, \Omega_c h^2, H_0, \sigma_8\}$, де $\sigma_8$~--- дисперсія флуктуацій густини матерії на масштабі $8\,h^{-1}$~Мпк (стандартний індикатор амплітуди структури на нелінійних масштабах), яка б одночасно задовольняла строгим обмеженням на кутовий спектр анізотропії CMB від місії \emph{Planck} та пояснювала надлишкову числову густину масивних галактик при $z \sim 10$--$12$ за даними космічного телескопа \emph{JWST}.
\end{enumerate}

%% --------------------------------------------------------
\section{Прогноз для майбутніх спостережень та верифікація}
%% --------------------------------------------------------

Якщо чисельний MCMC-аналіз підтвердить існування стійкої фізичної області параметрів для моделей із субгаббловою особливістю, це дасть змогу сформулювати конкретні передбачення. Вони можуть бути безпосередньо перевірені наступним поколінням астрофізичних інструментів у найближчі роки:

\begin{center}
\begin{tblr}{
  colspec = {llX},
  row{1}  = {font=\bfseries},
}
\toprule
Інструмент & Епоха даних & Об'єкт перевірки теоретичної моделі \\
\midrule
\emph{Euclid} \cite{euclid_2024}
  & Поточна / 2026+
  & Тривимірний спектр потужності $P(k)$ галактик на масштабах $k \sim 0.1$--$5\,h/\text{Мпк}$, що є критичним для точної фіксації положення зламу $k_\mathrm{break}$. \\
\emph{Nancy Grace Roman}
  & $\sim 2027+$
  & Огляди слабкого гравітаційного лінзування (weak lensing) та підрахунок галактик (galaxy counts) при $z \sim 2$--$4$, що накладе жорсткі обмеження на амплітуду нахилу $\Delta n$. \\
\emph{CMB-S4}
  & $\sim 2029+$
  & Надвисокі кутові мультиполі CMB ($\ell \sim 3000$--$5000$). Дослідження цієї області дозволить виявити тонкі ефекти в зоні затухання Сілка, викликані модифікацією первинного спектра. \\
\emph{SKA} (21~cm)
  & $\sim 2030+$
  & Радіоінтерферометричне картографування лінії 21~см у часи Темних віків. Це дасть прямий, нелінійний зріз спектра флуктуацій густини задовго до утворення перших зірок. \\
\bottomrule
\end{tblr}
\end{center}

Таким чином, розв'язання проблеми космологічних аномалій ранніх галактик лежить на перетині мікрофізики інфляційних потенціалів (таких як режими Ultra-Slow-Roll) та прецизійної великомасштабної астрофізики наступного десятиліття.
